\documentclass{beamer}
\usepackage{minted}
\usepackage{epigraph}
\usepackage{subfig}
\usepackage{tikz}
\usepackage{amsmath}
\usetheme{focus}

\title{Shor's Algorithm}
\subtitle{Quantum Computer Programming}
\author{Davide Carnemolla}
\titlegraphic{\includegraphics[scale=0.13]{images/logo.pdf}}
\institute{Department of Mathematics and \\ Computer Science \\ \\ University of Catania}
\date{2022/2023}

\begin{document}
    \begin{frame}
        \maketitle
    \end{frame}
    
    \section{Integer Factorization}

    \begin{frame}{The problem}
        \centering
        \includegraphics[height=2cm,keepaspectratio]{images/problem.pdf}
        \vspace{1cm}
        \begin{block}{Definition}
            Given an integer $N$, the goal of the \textit{integer factorization problem} is to find two primes, $p$ and $q$, such that $N = p \cdot q$.
        \end{block}
        \begin{exampleblock}{Example}
            Given $N = 143$ we know that the solution is the pair $(13, 11).$
        \end{exampleblock}
        
    \end{frame}

    \begin{frame}{RSA}
        \centering
        \begin{block}{RSA}
            $(pk,sk) \leftarrow$ \textbf{keygen}(): \hspace{0.05cm} given two primes $p$ and $q$ \\ 
                \hspace{3.7cm} let $n = pq$ \\ 
                \hspace{3.7cm} choose $e,d$ such that $e \cdot d \equiv 1 \; \text{mod} \; \phi(n)$ \\
                \hspace{3.7cm} \textbf{return} $pk = (n,e), \, sk = (n, d)$ \\
            \vspace{0.5cm}
            $c \leftarrow \textbf{enc}_{pk}(m) = m^e \; mod \; n$ \\
            \vspace{0.5cm}
            $m \leftarrow \textbf{dec}_{sk}(m) = c^d \; mod \; n$
        \end{block}
        \begin{alertblock}{RSA's Security}
            The RSA's Problem can be reduced to the Integer Factorization's Problem.
        \end{alertblock}
    \end{frame}

    \section{Solution in the classical model}
    \begin{frame}{Brute Force}
        \begin{alertblock}{Brute Force strategy}
            The brute force algorithm goes through all primes $p$ up to $\sqrt{N}$ and checks whether $p$ divides $N$.
        \end{alertblock}
        \begin{exampleblock}{Complexity}
            In the worst case, this would take time roughly , which is exponential in the number of digits $d = \log_2N$.
        \end{exampleblock}
    \end{frame}
    \begin{frame}{Quadratic Sieve}
        \begin{alertblock}{Quadratic Sieve Algorithm}
            A more efficient algorithm, known as the quadratic sieve, attempts to construct integers $a,b$ such that $a^2-b^2$ is a multiple of $N$.
            Once such $a,b$ are found, one checks whether $a b$ have common factors with $N$.
        \end{alertblock}
        \begin{exampleblock}{Complexity}
            The quadratic sieve method has asymptotic runtime exponential in $\sqrt{d}$, where $d = \log_2N$ is the number of digits of $N$.
        \end{exampleblock}
    \end{frame}
    \begin{frame}{General Number Field Sieve}
        \begin{alertblock}{GNFS}
            The Generic Number Field Sieve is most efficient classical factoring algorithm. The idea 
        \end{alertblock}
        \begin{exampleblock}{GNFS - Complexity}
            The number of digit of N is equals to $d = \log_2N$.
            The algorithm's complexity can be simplified as follows
            $$\mathcal{O}(exp(const \times d^{1/3}))$$
        \end{exampleblock}
    \end{frame}

    \section{The Period Finding Problem}
    \begin{frame}{The period finding problem}
        \begin{alertblock}{Definition of the problem}
            Given integers $N$ and $a$, find the smallest positive integer $r$ such that
            $$ a^r \equiv 1 \; mod \; N$$
            The number $r$ is called the period of $a$ modulo $N$.
        \end{alertblock}

        \begin{exampleblock}{Example}
            Suppose $N = 15$ and $a=7$ then
            \begin{equation*}
                    7^2 \equiv 4 \; mod \; 15, \quad
                    7^3 \equiv 13 \; mod \; 15, \quad
                    7^4 \equiv 1 \; mod \; 15 
            \end{equation*}
            That is, $7$ has period $4$ modulo $15$.
        \end{exampleblock}
    \end{frame}
    

    \begin{frame}{From factoring to period finding}
        \begin{columns}[t, onlytextwidth]
            \column{0.2\textwidth}
            \centering
            \Huge $(N, a)$

            \column{0.2\textwidth}
            \centering
            \Huge $\rightarrow$

            \column{0.2\textwidth}
            \centering
            \includegraphics[height=2.5cm,keepaspectratio]{images/robot.pdf}
            \vfill
            \footnotesize \textbf{period-finding machine}

            \column{0.2\textwidth}
            \centering
            \Huge $\rightarrow$

            \column{0.2\textwidth}
            \centering
            \Huge $r$ 
        \end{columns}
        \vspace{1.5cm}
        where $r$ is the period of $a$ modulo $N$.
    \end{frame}

    \section{Shor's Algorithm}

    \section{Compare between Shor and GNFS}

    \section{Implementation}
    
    
\end{document}
